\section{Security Incident: de la crisis a la mejora verificable}

Un identity provider es un objetivo de alto valor, y un incidente de seguridad desafía directamente los supuestos de confianza más fundamentales del cliente. La clase enfatiza el dolor que trae el incidente y la transformación cultural; el RCA oficial ofrece una perspectiva de sistema más concreta: cómo se expuso la credential, por qué la consulta de registros pasó eventos por alto, cómo se aprovechó el session token, cuándo se hizo la containment, cómo se notificó a los clientes y qué control se modificaron de forma permanente.

\subsection{Las cinco etapas de la incident response}

Un incident lifecycle completo incluye, como mínimo, detection, containment, investigation, remediation y notification. La detection puede provenir de la telemetry interna o de clientes y socios; la containment debe deshabilitar la account, hacer revoke de session/token y restringir las interfaces; la investigation establece el timeline y el affected scope; la remediation corrige credential, logging, workflow y product control; la notification, por su parte, entrega a los clientes afectados información sobre la que puedan actuar.

\lecturefigure{08-incident-response.png}{Un incidente de seguridad solo se da por terminado cuando se cierra el ciclo de aprendizaje de detection, containment, investigation, remediation y notification.}{Subtítulos locales 15:20--19:10; RCA del support-system de Okta de 2023; rediseño conceptual.}

\begin{knowledgebox}{Lectura de la figura: el RCA debe explicar «por qué no se detectó antes»}
La root cause no solo busca el error inicial; también debe explicar por qué las defensas no lo impidieron, por qué la telemetry no lo mostró, por qué la investigation query lo pasó por alto, por qué se retrasó el session revoke y cómo entraron las señales de los clientes en la respuesta. El RCA oficial de 2023 explica que distintas rutas de acceso a archivos generan distintos log event, por lo que la consulta inicial quedó incompleta; este detalle muestra que el propio esquema de observabilidad también forma parte de los security control.
\end{knowledgebox}

\begin{importantbox}{La containment y la remediation son distintas}
Deshabilitar la service account, hacer revoke de la session y bloquear direcciones IP es containment: su objetivo es detener el daño en curso. Impedir que los profile personales guarden credential de la empresa, reforzar la telemetry de acceso a archivos, vincular las session de administrador y mejorar la customer notification es remediation: su objetivo es reducir la probabilidad de recurrencia y acortar el detection time futuro. Ambas requieren distintos owner y distintas evidencias de cumplimiento.
\end{importantbox}

\begin{warningbox}{La transparencia no es publicar una entrada de blog}
La customer trust requiere un scope oportuno, los objetos afectados, la recommended action, lo que se sabe que se desconoce y actualizaciones posteriores. Dar demasiado pronto conclusiones definitivas no comprobadas se vuelve contra la confianza, y divulgar demasiado tarde le quita al cliente tiempo de containment. La incident communication debe ir a la par de la forensic confidence y conservar un decision log de quién sabía qué y en qué momento.
\end{warningbox}

\teachervoice{McKinnon dice que un incidente de seguridad es una de las etapas más difíciles de una carrera profesional, porque el headline asocia directamente a la empresa de identidad con «fue vulnerada». El juicio de la clase que vale la pena conservar es este: la confianza no se recupera con consignas; solo puede recuperarse convirtiendo el dolor en recursos, conductas y controles verificables orientados a security-first.}

\subsection{La security-first culture es un operating system}

La culture no es una value list colgada en la pared, sino lo que los líderes preguntan una y otra vez, lo que están dispuestos a retrasar en un launch, a quién asignan presupuesto, qué recompensan los ascensos y si las malas noticias se castigan. Si «security first» solo exige que el security team trabaje horas extra, mientras el product roadmap, los sales commitment y la executive review no cambian, la señal real que recibe la organización sigue siendo feature first.

\lecturefigure{09-security-first-culture.png}{La security-first culture se logra con la conducta de los líderes, los incentivos, los recursos y la evidencia en conjunto.}{Subtítulos locales 18:00--21:30; rediseño conceptual.}

\begin{knowledgebox}{Lectura de la figura: la culture puede auditarse}
En la leader behavior se observa si la alta dirección pregunta por los riesgos por iniciativa propia y acepta la escalation; en el incentive, los objetivos, el desempeño y los launch criteria; en el resource, si la renovación de la plataforma, el security staffing y el customer support reciben presupuesto; en la evidence, si los control test, los incident metric, las audit y los postmortem entran en la review. Aunque la culture es difícil de cuantificar con precisión, puede observarse a través de conductas repetidas y de la resource allocation.
\end{knowledgebox}

\subsection{Resumen de la sección}

El objetivo de la incident response no es «volver a la luz verde», sino establecer un learning loop que pueda revisarse. Security-first solo surte efecto en la siguiente situación de presión si se incorpora a las decisiones, a los recursos y a la evidencia de ingeniería.
